% !TeX root = ./ausarbeitung.tex

\section{Linear Least Squares Fit Data}

The functions given in the Analysis as \autoref{eq:least_squares:fit_polynomial} and \autoref{eq:least_squares:linear_fit} and reproduced here were fitted to the following thirteen independent measurements of $(x_i, y_i)$.

\begin{table}[htbp]
    \centering
    \caption{Given 13 Data Points for Applying Linear Least Squares Fit}
    \label{tab:appendix:least_squares_data}
    \(
    \begin{array}{
            c@{}
            S[table-format=+1.2]@{}
            c
            S[table-format=+1.1(2)]@{}
            c
        }
        ( &  {x}  & , &  {y}         & ) \\ \hline
        ( & -1.5  & , & -8.0 \pm 2.5 & ) \\
        ( & -1.25 & , & -2.5 \pm 2.0 & ) \\
        ( & -1.0  & , & -1.5 \pm 1.5 & ) \\
        ( & -0.75 & , &  1.0 \pm 1.0 & ) \\
        ( & -0.5  & , &  4.0 \pm 1.0 & ) \\
        ( & -0.25 & , &  4.0 \pm 0.5 & ) \\
        ( &  0.0  & , &  5.5 \pm 0.5 & ) \\
        ( & 0.25  & , &  4.5 \pm 0.5 & ) \\
        ( & 0.5   & , &  5.5 \pm 0.5 & ) \\
        ( & 0.75  & , &  3.0 \pm 1.0 & ) \\
        ( & 1.0   & , &  3.0 \pm 1.0 & ) \\
        ( & 1.25  & , &  4.5 \pm 1.5 & ) \\
        ( & 1.5   & , &  7.0 \pm 2.0 & )
    \end{array}
    \)
\end{table}

% applied fitfunctions
% \begin{gather}
%     y(x) = \theta_0 + \theta_1 x + \theta_2 x^2 + \theta_3 x^3
%     \label{eq:appendix:fit_polynomial} \\
%     y(x) = \theta_0 + \theta_1 x
%     \label{eq:appendix:linear_fit}
% \end{gather}


